\section{Full-Policy Verification Details}\label{sec:r23proofs}
This section supplies the arithmetic and implementation accounts for
Section~\ref{sec:r23policy}. All numerical coefficients are the exact values
of their stored binary64 representations. The stochastic law remains
Gaussian with those coefficients; no finite shock catalogue is substituted.

\subsection{The quadratic local advantage}
For a fixed supplied future policy, substitute
$U_{t+1}(z)=z'P_{t+1}z+c_{t+1}$ into its one-step cost. The action-dependent
terms are
\[
 u'H_tu+2\beta u'B_t'P_{t+1}A_tx.
\]
Since $H_t\succ0$, the minimum over the entire space $\R^d$ occurs at
$u_t^g=-H_t^{-1}\beta B_t'P_{t+1}A_tx$. Evaluating at $-K_tx$ and subtracting
this minimum yields
\[
 (-K_tx-u_t^g)'H_t(-K_tx-u_t^g)=x'D_t'H_t^{-1}D_tx.
\]
There is no approximation in this identity and no reference to the optimal
future policy. The possibly loose Frobenius bound in
\eqref{eq:r23eta} is chosen to make the implementation independent of an
uncertified spectral decomposition. A failed bound need not indicate that
the actual policy loss exceeds the tolerance.

\subsection{Outward matrix balls}
A matrix ball $(C,R)$ denotes every matrix $M$ with
$|M-C|\leq R$ entrywise, where $R\geq0$. Exact input balls have zero radius.
For conformable balls $(C,R)$ and $(D,S)$, let $k$ be the inner dimension.
The exact product is enclosed about the computed product $\operatorname{fl}(CD)$
by the radius
\begin{equation}
 \gamma_{2k}|C||D|+|C|S+R|D|+RS+2k\tau\mathbf1,
 \label{eq:r23ballproduct}
\end{equation}
with every nonnegative calculation rounded upward. Here
$\gamma_n=nu/(1-nu)$, $u$ is the binary64 unit roundoff, and $\tau$ is the
smallest positive normal binary64 number. The underflow allowance deliberately
exceeds the corresponding subnormal absolute-error allowance. The factor
$2k$ permits separate multiplication and addition rounding in classical
binary64 products. Fused operations only reduce this particular account.
Overflow and nonfinite values are rejected rather than certified.

To calculate a nonnegative product upper bound, the implementation divides
the computed product by a downward bound for $1-\gamma_{2k}$, adds the
underflow allowance, and rounds upward. Addition, scalar multiplication,
traces, squared Frobenius norms, and maximum absolute row sums have analogous
outward bounds. These operations enclose the actual recursion, not merely
a rounded residual evaluated without error bars. Positivity of $P_t$ follows
from the exact evaluation equation and positive stage costs, not from an
unverified numerical eigenvalue. In the executed model the positive diagonal
entries of $Q_t,R_t$ give the coercivity lower bounds directly. Covariance
validity is checked by symmetry and strict diagonal dominance.

The tests compare matrix operations with rational arithmetic on the exact
binary input values, including signed products and cancellation. Separate
tests compare policy losses with a separately computed optimal comparator,
check the coercive bound at states far outside an ordinary training range,
perturb policies within the declared execution class, and verify that the
certificate still runs when the Riccati comparator is disabled. These tests
complement, rather than replace, the enclosure argument.

\subsection{A robust action-implementation envelope}
Suppose the implemented action is $-K_tx+e_t$, with
$\|e_t\|_2\leq e\|x\|_2$, possibly depending on the past and current state.
Put $\overline P_T=Q_T$ and $\overline c_T=0$. For $F_t=A_t-B_tK_t$, set
\begin{align}
 h_t={}&2\|R_tK_t\|_F e+\|R_t\|_\infty e^2\nonumber\\
 &+\beta\{2\|\overline P_{t+1}F_t\|_F\|B_t\|_F e
             +\|\overline P_{t+1}\|_\infty\|B_t\|_F^2e^2\},
 \label{eq:r23implh}\\
 \overline P_t={}&Q_t+K_t'R_tK_t+
                    \beta F_t'\overline P_{t+1}F_t+h_tI,\nonumber\\
 \overline c_t={}&\beta\{\overline c_{t+1}+\tr(\overline P_{t+1}\Sigma)\}.
 \label{eq:r23implrec}
\end{align}
All displayed coefficients are replaced by their outward upper enclosures
in the numerical calculation.

\begin{proposition}[Implementation allowance]\label{prop:r23implementation}
For every execution satisfying the relative action-error bound, its remaining
cost is at most $x'\overline P_tx+\overline c_t$. On the initial ball
$\|x\|_2^2\leq b$, its excess over the ideal stored policy is at most
\begin{equation}
 \varepsilon_{\rm impl}
 =\max\{0,b\|\overline P_0-P_0\|_\infty+
                         \overline c_0-c_0\}.
 \label{eq:r23implallow}
\end{equation}
\end{proposition}
\begin{proof}
Expanding the current action cost around $-K_tx$ bounds its increase by
$\{2\|R_tK_t\|_F e+\|R_t\|_\infty e^2\}\|x\|_2^2$.
Expanding the next-state quadratic envelope around $F_tx$ bounds its
increase by the second line of \eqref{eq:r23implh}, times $\|x\|_2^2$.
The mean-zero innovation contributes its trace term and has no cross term
with a nonanticipating action. Backward induction proves the value envelope.
Subtract the exact stored-policy quadratic value, use symmetry and the
maximum-row-sum bound, and maximize over the initial ball. This proves
\eqref{eq:r23implallow} even when the execution errors are history dependent.
\end{proof}

The error budget is explicit. Certification of arbitrary machine executions
without an error contract would be a different claim. The experiment neither
removes this term nor infers it from agreement on sampled states.

\subsection{Reoptimized comparative statics}
Let $K^{(r)}$ be a certified complete policy in regime $r$, and let its
reported total loss bound on the initial ball be $b_r$. At a fixed initial
state, the first-action objective with the \emph{optimal} future is a strictly
convex quadratic with Hessian at least $2R_0$. The cost difference between
using the candidate's first action and the optimal first action, both followed
by the optimal future, is no larger than the candidate's full policy loss.
Thus, with $r_0\leq\lambda_{\min}(R_0)$,
\begin{equation}
 \|u_0^{(r)}-u_0^{*,(r)}\|_2\leq\sqrt{b_r/r_0}.
 \label{eq:r23actionradius}
\end{equation}
At $x_0=-\mathbf1$, define mean investment
$m_r=d^{-1}\mathbf1'K_0^{(r)}\mathbf1$. Its distance from the optimal mean
is at most $\sqrt{b_r/(r_0d)}$. The reported radius also includes the
relative-execution and summation allowances. Subtracting the interval for
the anchor from that for the changed regime gives a valid optimum-response
interval. It is an enclosure in a specified model, not a confidence interval
for an estimated economic population. Strong competitors are evaluated by
the same construction.

\subsection{Trainable continuation factors}
For the square-activation critic, write
$L(W;P)=\frac14\|W'W-dP\|_F^2$, with symmetric own-policy target $P$.
Differentiation gives $\nabla_W L=W(W'W-dP)$. If $X\sim N(0,I)$, the mean
squared error between the gradients of $\|WX\|^2/d$ and $X'PX$ equals
$4\|W'W-dP\|_F^2/d^2$. The target is obtained by evaluating the current
candidate policy; it is not the Riccati value matrix. The experiment trains
all hidden weights for a fixed number of gradient updates and then solves
the actor's full vector-action problem. A conventional quadratic receives
the identical population target and is allowed its exact minimizer. This
matched-information comparison isolates the cost of the neural factorization.
No advantage over conventional quadratics is built into the target oracle.
